\documentclass[../../main.tex]{subfiles}% 注意这里的文件路径不能用 ./main.tex ,否则用latexmk编译子文件会报错
\graphicspath{{\subfix{./image/}}{\subfix{../../image/}}} % 指定图片引用路径,后续可以直接使用图片文件名
% 如果图片存放在上级目录,则使用 ../../image/ 作为备用路径

% 例如:
% \begin{figure}[H]
% \centering
% \includegraphics[width=0.3\textwidth]{图.png}
% \caption{}
% \label{figure:图}
% \end{figure}
% 注意:上述\label{}一定要放在\caption{}之后,否则引用图片序号会只会显示??.

\begin{document}

\section{存在性结果的证明}\label{section:ch07-存在性证明}

\cref{theorem:ch07-7.3.1}的证明依赖下列结果，它是\cref{theorem:ch06-6.6.1}的对应物。如\cref{theorem:ch06-6.6.1}情形，证明需要若干预备结果。

\begin{theorem}\label{theorem:ch07-7.5.1}
设 (A0)--(A1) 成立，$\{u_n\}$ 为问题 (Pn) 的熵解序列。则存在子列 $\{u_{n_k}\}\subseteq\{u_n\}$，具有下列性质：
\begin{enumerate}[(i)]
\item 存在 $u\in L^\infty_{w^*}(0,T;\mathscr{R}_f^+(\mathbb{R}))$ 使
\begin{align}
u_{n_k}\stackrel{*}{\rightharpoonup}u\quad\text{在 }L^\infty_{w^*}(0,T;\mathscr{R}_f(\mathbb{R})); \label{align:ch07-7.88}
\end{align}
\item 与 $\{u_{n_k}\}$ 相伴的 Young 测度序列 $\{\nu_{n_k}\}\subseteq\mathcal{Y}^+(S;\mathbb{R})$ 窄收敛于 Young 测度 $\nu\in\mathcal{Y}^+(S;\mathbb{R})$，且对几乎所有 $(x,t)\in S$，
\begin{align}
u_r(x,t)=\int_{[0,\infty)}y\,\mathrm{d}\nu_{(x,t)}(y), \label{align:ch07-7.89}
\end{align}
\begin{align}
\phi(u_r(x,t))=\int_{[0,\infty)}\phi(y)\,\mathrm{d}\nu_{(x,t)}(y), \label{align:ch07-7.90}
\end{align}
其中 $\{\nu_{(x,t)}\}_{(x,t)\in S}$ 为 $\nu$ 的离散分解；
\item 对每个满足
\begin{align}
\lim_{y\to\infty}\frac{U(y)}{y}=:M_U\in\mathbb{R} \label{align:ch07-7.91}
\end{align}
的 $U\in C([0,\infty))$ 与每个 $\zeta\in C([0,T];C_c(\mathbb{R}))$，有
\begin{align}
\lim_{k\to\infty}\iint_S U(u_{n_k})\zeta\,\mathrm{d}x\,\mathrm{d}t
=\iint_S U^*\zeta\,\mathrm{d}x\,\mathrm{d}t+M_U\int_0^T\langle u_s(\cdot,t),\zeta(\cdot,t)\rangle\,\mathrm{d}t, \label{align:ch07-7.92}
\end{align}
其中 $U^*\in L^\infty(0,T;L^1_{\mathrm{loc}}(\mathbb{R}))$ 定义为
\begin{align}
U^*(x,t)\triangleq \int_{[0,\infty)}U(y)\,\mathrm{d}\nu_{(x,t)}(y)
\quad\text{对几乎所有 }(x,t)\in S. \label{align:ch07-7.93}
\end{align}
\end{enumerate}
\end{theorem}

为证\cref{theorem:ch07-7.5.1}，需要下列结果。

\begin{proposition}\label{proposition:ch07-7.5.2}
设 (A0)--(A1) 成立，$\{u_n\}$ 为问题 (Pn) 的熵解序列。则存在子列 $\{u_{n_k}\}\subseteq\{u_n\}$、Young 测度 $\nu\in\mathcal{Y}^+(S;\mathbb{R})$ 与 $\sigma\in L^\infty_{w^*}(0,T;\mathscr{R}_f^+(\mathbb{R}))$，使
\begin{align}
u_{n_k}\stackrel{*}{\rightharpoonup}u_b+\sigma\quad\text{在 }L^\infty_{w^*}(0,T;\mathscr{R}_f(\mathbb{R})), \label{align:ch07-7.94}
\end{align}
其中 $u_b\in L^\infty(0,T;L^1(\mathbb{R}))$ 定义为
\begin{align}
u_b(x,t)\triangleq \int_{[0,\infty)}y\,\mathrm{d}\nu_{(x,t)}(y)
\quad\text{对几乎所有 }(x,t)\in S. \label{align:ch07-7.95}
\end{align}
\end{proposition}

\begin{proof}
由\eqref{align:ch07-7.71}中第二个不等式，序列 $\{u_n\}$ 在 $L^\infty(0,T;L^1(\mathbb{R}))\subseteq L^\infty_{w^*}(0,T;\mathscr{R}_f(\mathbb{R}))$ 中有界。由\cref{proposition:ch04-4.4.16}与 Banach--Alaoglu 定理，存在 $\mu\in L^\infty_{w^*}(0,T;\mathscr{R}_f^+(\mathbb{R}))$ 使
\begin{align}
u_{n_k}\stackrel{*}{\rightharpoonup}\mu\quad\text{在 }L^\infty_{w^*}(0,T;\mathscr{R}_f(\mathbb{R})). \label{align:ch07-7.96}
\end{align}
另一方面，因 $\{u_{n_k}\}$ 在 $L^1(S)$ 中有界，由\cref{remark:ch05-5.4.4}(ii)，存在 $\{u_{n_k}\}$ 的子列（不重新标号）与 $\sigma\in\mathscr{R}_f^+(S)$ 使
\begin{align}
u_{n_k}\stackrel{*}{\rightharpoonup}u_b+\sigma\quad\text{在 }\mathscr{R}_f(S). \label{align:ch07-7.97}
\end{align}
由\eqref{align:ch07-7.96}--\eqref{align:ch07-7.97}得 $\mu=u_b+\sigma$，结论得证。
\end{proof}

\begin{proposition}\label{proposition:ch07-7.5.3}
设 (A0)--(A1) 成立，$\{u_n\}$ 为问题 (Pn) 的熵解序列。设 $\{u_{n_k}\}\subseteq\{u_n\}$、$\nu\in\mathcal{Y}^+(S;\mathbb{R})$ 与 $\sigma\in L^\infty_{w^*}(0,T;\mathscr{R}_f^+(\mathbb{R}))$ 由\cref{proposition:ch07-7.5.2}给出。设 $U\in C([0,\infty))$ 满足\eqref{align:ch07-7.91}。则对任意 $\zeta\in C([0,T];C_c(\mathbb{R}))$，
\begin{align}
\lim_{k\to\infty}\iint_S U(u_{n_k})\zeta\,\mathrm{d}x\,\mathrm{d}t
=\iint_S U^*\zeta\,\mathrm{d}x\,\mathrm{d}t+M_U\int_0^T\langle\sigma(\cdot,t),\zeta(\cdot,t)\rangle\,\mathrm{d}t \label{align:ch07-7.98}
\end{align}
其中 $U^*\in L^\infty(0,T;L^1_{\mathrm{loc}}(\mathbb{R}))$ 由\eqref{align:ch07-7.93}给出。
\end{proposition}

\begin{proof}
固定任意 $L>0$。由\eqref{align:ch07-7.91}，对任意 $\epsilon>0$，存在 $y_\epsilon>0$ 使对所有 $y>y_\epsilon$，
\begin{align}
-\epsilon y<U(y)-M_U y<\epsilon y. \label{align:ch07-7.99}
\end{align}
对任意 $j\in\mathbb{N}$、$j\geqslant[y_\epsilon]+1$，设 $f_{j1},f_{j2}\in C([0,\infty))$ 满足
\begin{align}
\begin{cases}
0\leqslant f_{j1}\leqslant1,\ 0\leqslant f_{j2}\leqslant1,\ f_{j1}+f_{j2}=1\text{ 在 }[0,\infty),\\
\operatorname{supp}f_{j1}\subseteq[0,j+1],\ \operatorname{supp}f_{j2}\subseteq[j,\infty).
\end{cases} \label{align:ch07-7.100}
\end{align}
利用\eqref{align:ch07-7.99}与\eqref{align:ch07-7.71}，对固定的 $k\in\mathbb{N}$ 与 $t\in(0,T)$，
\begin{align*}
\begin{aligned}
\int_{I_L}|U(u_{n_k}(x,t))|\,\mathrm{d}x
&=\int_{I_L}f_{j1}(u_{n_k}(x,t))|U(u_{n_k}(x,t))|\,\mathrm{d}x+\int_{I_L}f_{j2}(u_{n_k}(x,t))|U(u_{n_k}(x,t))|\,\mathrm{d}x\\
&\leqslant2L\|U\|_{L^\infty([0,j+1])}+\int_{\{u_{n_k}(\cdot,t)>y_\epsilon\}}|U(u_{n_k}(x,t))|\,\mathrm{d}x\\
&\leqslant2L\|U\|_{L^\infty([0,j+1])}+(|M_U|+\epsilon)\int_{I_L}|u_{n_k}(x,t)|\,\mathrm{d}x\\
&\leqslant2L\|U\|_{L^\infty([0,j+1])}+(|M_U|+\epsilon)\|u_0\|_{\mathscr{R}_f(\mathbb{R})}
\end{aligned}
\end{align*}
（回忆 $I_L=(-L,L)$），故
\begin{align}
\|U(u_{n_k})\|_{L^\infty(0,T;L^1(I_L))}\leqslant2L\|U\|_{L^\infty([0,j+1])}+(|M_U|+\epsilon)\|u_0\|_{\mathscr{R}_f(\mathbb{R})}. \label{align:ch07-7.101}
\end{align}
由 Banach--Alaoglu 定理，存在 $\{U(u_{n_k})\}$ 的子列（不重新标号）与 $\mu\in L^\infty_{w^*}(0,T;\mathscr{R}_f(I_L))$ 使
\begin{align}
U(u_{n_k})\stackrel{*}{\rightharpoonup}\mu\quad\text{在 }L^\infty_{w^*}(0,T;\mathscr{R}_f(I_L)). \label{align:ch07-7.102}
\end{align}
此外，如\cref{proposition:ch05-5.4.13}证明中那样论证，
\begin{align}
U(u_{n_k})\stackrel{*}{\rightharpoonup}U^*+M_U\sigma\quad\text{在 }\mathscr{R}_f(R_L). \label{align:ch07-7.103}
\end{align}
由\eqref{align:ch07-7.102}--\eqref{align:ch07-7.103}，对任意 $\zeta\in C([0,T];C_c(I_L))$，
\begin{align}
\lim_{k\to\infty}\iint_{R_L}U(u_{n_k})\zeta\,\mathrm{d}x\,\mathrm{d}t
=\iint_{R_L}U^*\zeta\,\mathrm{d}x\,\mathrm{d}t+M_U\int_0^T\langle\sigma(\cdot,t),\zeta(\cdot,t)\rangle\,\mathrm{d}t \label{align:ch07-7.104}
\end{align}
其中 $U^*$ 由\eqref{align:ch07-7.93}在 $R_L$ 中几乎处处定义。

等式\eqref{align:ch07-7.98}由\eqref{align:ch07-7.104}得到：因对任意 $\zeta\in C([0,T];C_c(\mathbb{R}))$，可取 $L>0$ 充分大使 $\operatorname{supp}\zeta\subseteq((-L,L)\times[0,T])$。此外，由\eqref{align:ch07-7.91}存在 $M>0$ 使 $|U(y)|\leqslant M(y+1)$（$y\geqslant0$），故 $|U^*|\leqslant M(u_b+1)$ 在 $S$ 中几乎处处（见\eqref{align:ch07-7.93}与\eqref{align:ch07-7.95}）。因 $u_b\in L^\infty(0,T;L^1(\mathbb{R}))$，得 $U^*\in L^\infty(0,T;L^1_{\mathrm{loc}}(\mathbb{R}))$。
\end{proof}

为继续，需要下列引理，其证明在本节末给出。

\begin{lemma}\label{lemma:ch07-7.5.4}
设 (A0)--(A1) 成立，$U\in C([0,\infty))$ 满足\eqref{align:ch07-7.91}。设 $\{u_{n_k}\}$、$\sigma$ 与 $U^*\in L^\infty(0,T;L^1_{\mathrm{loc}}(\mathbb{R}))$ 如\cref{proposition:ch07-7.5.3}。则对每个 $\zeta\in C([0,T];C_c(\mathbb{R}))$，
\begin{align}
\lim_{k\to\infty}\int_0^T\left|\int_{\mathbb{R}}U(u_{n_k}(x,t))\zeta(x,t)\,\mathrm{d}x
-\int_{\mathbb{R}}U^*(x,t)\zeta(x,t)\,\mathrm{d}x
-M_U\langle\sigma(\cdot,t),\zeta(\cdot,t)\rangle\right|\mathrm{d}t=0. \label{align:ch07-7.105}
\end{align}
\end{lemma}

\begin{remark}\label{remark:ch07-7.5.1}
取 $\zeta(x,t)=\rho(x)$、$\rho\in C_c(\mathbb{R})$，由\eqref{align:ch07-7.105}得
\begin{align}
\lim_{k\to\infty}\int_0^T\left|\int_{\mathbb{R}}U(u_{n_k}(x,t))\rho(x)\,\mathrm{d}x-U_\rho^*(t)-M_U\langle\sigma(\cdot,t),\rho\rangle\right|\mathrm{d}t=0, \label{align:ch07-7.106a}
\end{align}
其中
\begin{align}
U_\rho^*(t)\triangleq \int_{\mathbb{R}}U^*(x,t)\rho(x)\,\mathrm{d}x. \label{align:ch07-7.106b}
\end{align}
由\eqref{align:ch07-7.106a}与 $C_c(\mathbb{R})$ 的可分性，对每个 $L>0$ 与如\cref{lemma:ch07-7.5.4}的 $U$，存在零测集 $N\subseteq(0,T)$ 与 $\{u_{n_k}\}$ 的子列（不重新标号），使对所有 $t\in(0,T)\setminus N$，
\begin{align}
U(u_{n_k})(\cdot,t)\stackrel{*}{\rightharpoonup}U^*(\cdot,t)+M_U\sigma(\cdot,t)\quad\text{在 }\mathscr{R}_f(I_L). \label{align:ch07-7.107}
\end{align}
\end{remark}

\begin{proposition}\label{proposition:ch07-7.5.5}
设 (A0)--(A1) 成立，$\{u_n\}$ 为问题 (Pn) 的熵解序列。设 $\{u_{n_k}\}\subseteq\{u_n\}$ 与 $\sigma\in L^\infty_{w^*}(0,T;\mathscr{R}_f^+(\mathbb{R}))$ 由\cref{proposition:ch07-7.5.2}给出，并按
\begin{align}
u\triangleq u_b+\sigma \label{align:ch07-7.108}
\end{align}
定义 $u\in L^\infty_{w^*}(0,T;\mathscr{R}_f^+(\mathbb{R}))$。则
\begin{align}
u_r=u_b\quad\text{在 }S\text{ 中几乎处处}, \label{align:ch07-7.109a}
\end{align}
\begin{align}
u_s=\sigma\quad\text{在 }L^\infty_{w^*}(0,T;\mathscr{R}_f(\mathbb{R})). \label{align:ch07-7.109b}
\end{align}
\end{proposition}

\begin{proof}
对任意非负 $\zeta\in C^1([0,T];C^1_c(\mathbb{R}))$、$\zeta(\cdot,T)=0$，将\eqref{align:ch07-7.77}与\eqref{align:ch07-7.78}相加，得
\begin{align*}
\iint_S\{U_j(u_{n_k})\partial_t\zeta+\Lambda_{U_j}(u_{n_k})\partial_x\zeta\}\,\mathrm{d}x\,\mathrm{d}t
\geqslant-\int_{\mathbb{R}}U_j(u_{0n_k})(x)\zeta(x,0)\,\mathrm{d}x,
\end{align*}
其中对所有 $j\in\mathbb{N}$，
\begin{align}
U_j(y)\triangleq (y-j)\chi_{[j,\infty)}(y),\quad
\Lambda_{U_j}(y)\triangleq \int_0^y U_j'(z)\phi'(z)\,\mathrm{d}z=[\phi(y)-\phi(j)]\chi_{[j,\infty)}(y). \label{align:ch07-7.110}
\end{align}
由标准逼近论证，对上述每个 $\zeta$ 与几乎所有 $\tau\in(0,T)$，
\begin{align}
\begin{aligned}
\int_{\mathbb{R}}&U_j(u_{n_k})(x,\tau)\zeta(x,\tau)\,\mathrm{d}x-\int_{\mathbb{R}}U_j(u_{0n_k})(x)\zeta(x,0)\,\mathrm{d}x\\
&\leqslant\iint_{S_\tau}\{U_j(u_{n_k})\partial_t\zeta+\Lambda_{U_j}(u_{n_k})\partial_x\zeta\}\,\mathrm{d}x\,\mathrm{d}t,
\end{aligned} \label{align:ch07-7.111}
\end{align}
其中 $S_\tau\triangleq \mathbb{R}\times(0,\tau)$。因 $M_{U_j}=1$、$M_{\Lambda_{U_j}}=0$，由\eqref{align:ch07-7.105}得
\begin{align}
\lim_{k\to\infty}\int_0^\tau\left|\int_{\mathbb{R}}[U_j(u_{n_k})\partial_t\zeta](x,t)\,\mathrm{d}x
-\int_{\mathbb{R}}[U_j^*\partial_t\zeta](x,t)\,\mathrm{d}x
-\langle\sigma(\cdot,t),\partial_t\zeta(\cdot,t)\rangle\right|\mathrm{d}t=0, \label{align:ch07-7.112a}
\end{align}
\begin{align}
\lim_{k\to\infty}\int_0^\tau\left|\int_{\mathbb{R}}[\Lambda_{U_j}(u_{n_k})\partial_x\zeta](x,t)\,\mathrm{d}x
-\int_{\mathbb{R}}[\Lambda_{U_j}^*\partial_x\zeta](x,t)\,\mathrm{d}x\right|\mathrm{d}t=0, \label{align:ch07-7.112b}
\end{align}
其中
\begin{align*}
U_j^*\triangleq \int_{[0,\infty)}U_j(y)\,\mathrm{d}\nu(y),\qquad
\Lambda_{U_j}^*\triangleq \int_{[0,\infty)}\Lambda_{U_j}(y)\,\mathrm{d}\nu(y).
\end{align*}
由\eqref{align:ch07-7.112a}得
\begin{align}
\begin{aligned}
\lim_{k\to\infty}\iint_{S_\tau}&\{U_j(u_{n_k})\partial_t\zeta+\Lambda_{U_j}(u_{n_k})\partial_x\zeta\}\,\mathrm{d}x\,\mathrm{d}t\\
&=\iint_{S_\tau}\{U_j^*\partial_t\zeta+\Lambda_{U_j}^*\partial_x\zeta\}\,\mathrm{d}x\,\mathrm{d}t
+\int_0^\tau\langle\sigma(\cdot,t),\partial_t\zeta(\cdot,t)\rangle\,\mathrm{d}t.
\end{aligned} \label{align:ch07-7.113}
\end{align}
另一方面，由\eqref{align:ch07-7.107}与对角论证，存在零测集 $N\subseteq(0,T)$ 与 $\{u_{n_k}\}$ 的子列（不重新标号），使对所有 $\tau\in(0,T)\setminus N$ 与 $j\in\mathbb{N}$，
\begin{align}
\lim_{k\to\infty}\int_{\mathbb{R}}U_j(u_{n_k})(x,\tau)\zeta(x,\tau)\,\mathrm{d}x
=\int_{\mathbb{R}}U_j^*(x,\tau)\zeta(x,\tau)\,\mathrm{d}x+\langle\sigma(\cdot,\tau),\zeta(x,\tau)\rangle. \label{align:ch07-7.114}
\end{align}
此外，因序列 $\{U_j(u_{0n_k})-u_{0n_k}\}$ 在 $L^\infty(\mathbb{R})$ 中有界且
\begin{align*}
\lim_{k\to\infty}(U_j(u_{0n_k})-u_{0n_k})=U_j(u_{0r})-u_{0r}\ \text{在 }\mathbb{R}\text{ 中几乎处处}
\end{align*}
（见\eqref{align:ch07-7.45}），由\eqref{align:ch07-7.45}得
\begin{align}
\lim_{k\to\infty}\int_{\mathbb{R}}U_j(u_{0n_k})(x)\zeta(x,0)\,\mathrm{d}x
=\int_{\mathbb{R}}U_j(u_{0r})(x)\zeta(x,0)\,\mathrm{d}x+\langle u_{0s},\zeta(\cdot,0)\rangle. \label{align:ch07-7.115}
\end{align}
在\eqref{align:ch07-7.111}中令 $k\to\infty$，由\eqref{align:ch07-7.113}--\eqref{align:ch07-7.115}，对所有 $\tau\in(0,T)\setminus N$ 与 $j\in\mathbb{N}$，得
\begin{align}
\begin{aligned}
\int_{\mathbb{R}}U_j^*(x,\tau)\zeta(x,\tau)\,\mathrm{d}x+\langle\sigma(\cdot,\tau),\zeta(\cdot,\tau)\rangle
\leqslant{}&\iint_{S_\tau}\{U_j^*\partial_t\zeta+\Lambda_{U_j}^*\partial_x\zeta\}\,\mathrm{d}x\,\mathrm{d}t\\
&+\int_0^\tau\langle\sigma(\cdot,t),\partial_t\zeta(\cdot,t)\rangle\,\mathrm{d}t
+\int_{\mathbb{R}}U_j(u_{0r})(x)\zeta(x,0)\,\mathrm{d}x+\langle u_{0s},\zeta(\cdot,0)\rangle.
\end{aligned} \label{align:ch07-7.116}
\end{align}
因对所有 $y\geqslant0$，$0\leqslant U_j(y)\leqslant y\chi_{[j,\infty)}(y)$ 且 $|\Lambda_{U_j}(y)|\leqslant My\chi_{[j,\infty)}(y)$（见\eqref{align:ch07-7.1}），故 $|U_j^*|\leqslant u_b$ 且 $|\Lambda_{U_j}^*|\leqslant Mu_b$（$j\in\mathbb{N}$）。因 $u_b\in L^\infty(0,T;L^1(\mathbb{R}))$，显然
\begin{align}
\lim_{j\to\infty}U_j^*=\lim_{j\to\infty}\Lambda_{U_j}^*=0\quad\text{在 }S\text{ 中几乎处处且在 }L^1(S), \label{align:ch07-7.117}
\end{align}
从而在\eqref{align:ch07-7.116}中令 $j\to\infty$，对所有 $\tau\in(0,T)\setminus N$ 与 $\zeta\in C^1([0,T];C^1_c(\mathbb{R}))$ 得
\begin{align}
\langle\sigma(\cdot,\tau),\zeta(\cdot,\tau)\rangle
\leqslant\int_0^\tau\langle\sigma(\cdot,t),\partial_t\zeta(\cdot,t)\rangle\,\mathrm{d}t+\langle u_{0s},\zeta(\cdot,0)\rangle. \label{align:ch07-7.118}
\end{align}
在\eqref{align:ch07-7.118}中取 $\zeta(x,t)=\rho(x)$，得对任意非负 $\rho\in C^1_c(\mathbb{R})$，$\langle\sigma(\cdot,\tau),\rho\rangle\leqslant\langle u_{0s},\rho\rangle$，故 $\sigma(\cdot,\tau)\ll u_{0s}$（回忆 $\sigma$ 非负）。于是对任意 $\tau\in(0,T)\setminus N$，$\sigma(\cdot,\tau)$ 关于 Lebesgue 测度奇异。另一方面，因 $u_b\in L^\infty(0,T;L^1(\mathbb{R}))$，测度 $E\mapsto\iint_E u_b\,\mathrm{d}x\,\mathrm{d}t$（$E\in\mathcal{B}(\mathbb{R}^2)\cap S$）关于 Lebesgue 测度绝对连续。由 Lebesgue 分解的唯一性（见\cref{theorem:ch01-Lebesgue分解}）得结论。
\end{proof}

\begin{proposition}\label{proposition:ch07-divcurl}\label{proposition:ch07-7.5.6}
设 (A0)--(A1) 成立。设 $\nu\in\mathcal{Y}^+(S;\mathbb{R})$ 由\cref{proposition:ch07-7.5.2}给出，$u\in L^\infty_{w^*}(0,T;\mathscr{R}_f(\mathbb{R}))$ 由\eqref{align:ch07-7.108}定义。则
\begin{align}
\phi^*(x,t)=\phi(u_r(x,t))\quad\text{对几乎所有 }(x,t)\in S \label{align:ch07-7.119}
\end{align}
其中 $\phi^*$ 由\eqref{align:ch07-7.4}定义。
\end{proposition}

\begin{proof}
设 $U,V\in C^2([0,\infty))\cap L^\infty(0,\infty)$ 满足\eqref{align:ch07-7.58}，并设 $\Lambda_U,\Lambda_V\in L^\infty(0,\infty)$（见\eqref{align:ch07-7.83}）。由\cref{lemma:ch07-7.4.8}，对任意 $n\in\mathbb{N}$，存在 $\sigma_n,\tau_n\in\mathscr{R}_f(S)$，使对所有 $\zeta\in C^1_c(S)$，
\begin{align}
\iint_S\{U(u_n)\partial_t\zeta+\Lambda_U(u_n)\partial_x\zeta\}\,\mathrm{d}x\,\mathrm{d}t=\langle\sigma_n,\zeta\rangle, \label{align:ch07-7.120a}
\end{align}
\begin{align}
\iint_S\{V(u_n)\partial_t\zeta+\Lambda_V(u_n)\partial_x\zeta\}\,\mathrm{d}x\,\mathrm{d}t=\langle\tau_n,\zeta\rangle, \label{align:ch07-7.120b}
\end{align}
其中 $u_n$ 为\cref{proposition:ch07-7.4.7}给出的 (Pn) 的熵解。

设 $A\Subset S$ 为有界开集，定义 $Y_n,Z_n:A\to\mathbb{R}^2$ 为
\begin{align*}
Y_n\triangleq (\Lambda_U(u_n),U(u_n)),\qquad Z_n\triangleq (V(u_n),-\Lambda_V(u_n)).
\end{align*}
由\eqref{align:ch07-7.120a}得
\begin{align}
\operatorname{div}Y_n=-\sigma_n,\qquad \operatorname{curl}Z_n=-\tau_n\quad\text{在 }\mathcal{D}^*(A). \label{align:ch07-7.121}
\end{align}
因 $U,\Lambda_U,V,\Lambda_V$ 在 $(0,\infty)$ 上有界，序列 $\{U(u_n)\}$、$\{\Lambda_U(u_n)\}$、$\{V(u_n)\}$、$\{\Lambda_V(u_n)\}$ 在 $L^\infty(A)$ 中有界，从而在 $L^1(A)$ 中有界且一致可积。由\cref{proposition:ch05-5.4.10}，在 $L^1(A)$ 中
\begin{align*}
U(u_n)\rightharpoonup U^*\triangleq \int_{[0,\infty)}U(y)\,\mathrm{d}\nu(y),\quad
\Lambda_U(u_n)\rightharpoonup\Lambda_U^*\triangleq \int_{[0,\infty)}\Lambda_U(y)\,\mathrm{d}\nu(y),
\end{align*}
\begin{align*}
V(u_n)\rightharpoonup V^*\triangleq \int_{[0,\infty)}V(y)\,\mathrm{d}\nu(y),\quad
\Lambda_V(u_n)\rightharpoonup\Lambda_V^*\triangleq \int_{[0,\infty)}\Lambda_V(y)\,\mathrm{d}\nu(y).
\end{align*}
同理，$\{U(u_n)\}$、$\{\Lambda_U(u_n)\}$、$\{V(u_n)\}$、$\{\Lambda_V(u_n)\}$ 有界，从而在 $L^2(A)$ 中弱收敛。于是当 $n\to\infty$ 时，
\begin{align}
Y_n\rightharpoonup Y^*\triangleq (\Lambda_U^*,U^*),\quad
Z_n\rightharpoonup Z^*\triangleq (V^*,-\Lambda_V^*)\quad\text{在 }L^2(A;\mathbb{R}^2). \label{align:ch07-7.122}
\end{align}
类似地，
\begin{align}
Y_n\cdot Z_n\triangleq \Lambda_U(u_n)V(u_n)-\Lambda_V(u_n)U(u_n)
\rightharpoonup\int_{[0,\infty)}[\Lambda_U(y)V(y)-\Lambda_V(y)U(y)]\,\mathrm{d}\nu(y)
\quad\text{在 }L^2(A). \label{align:ch07-7.123}
\end{align}
由\eqref{align:ch07-7.85}（应用于 $\sigma_n$ 与 $\tau_n$）与\eqref{align:ch07-7.121}，$\{\operatorname{div}Y_n\}$ 与 $\{\operatorname{curl}Z_n\}$ 在 $H^{-1}(A)$ 中预紧。于是由\eqref{align:ch07-7.122}与 div--curl 引理（Murat--Tartar）得
\begin{align}
Y_n\cdot Z_n\to Y^*\cdot Z^*=\Lambda_U^*V^*-\Lambda_V^*U^*\quad\text{在 }\mathcal{D}^*(A) \label{align:ch07-7.124}
\end{align}
（例如见文献~\cite{ref46} 推论 1.3.1 与定理 5.2.1）。由\eqref{align:ch07-7.123}--\eqref{align:ch07-7.124}得
\begin{align}
\int_{[0,\infty)}[\Lambda_U(y)-\Lambda_U^*]V(y)\,\mathrm{d}\nu(y)
=\int_{[0,\infty)}[U(y)-U^*]\Lambda_V(y)\,\mathrm{d}\nu(y)
\quad\text{在 }A\text{ 中几乎处处}. \label{align:ch07-7.125}
\end{align}

设 $U$ 如上，$U'>0$ 在 $(0,\infty)$。由标准逼近论证可取 $V(y)=|U(y)-U^*|$，适当选取 $\vartheta_V$ 得
\begin{align*}
\Lambda_V(y)=\int_0^y V'(z)\phi'(z)\,\mathrm{d}z+\vartheta_V
=\operatorname{sgn}(U(y)-U^*)[\Lambda_U(y)-\Lambda_U(U^{-1}(U^*))],
\end{align*}
故由\eqref{align:ch07-7.125}，
\begin{align}
[\Lambda_U^*-\Lambda_U(U^{-1}(U^*))]\int_{[0,\infty)}|U(y)-U^*|\,\mathrm{d}\nu(y)=0
\quad\text{在 }A\text{ 中几乎处处}. \label{align:ch07-7.126}
\end{align}
设 $U_p\in C^2([0,\infty))\cap L^\infty(0,\infty)$，$U_p(0)=0$（$p\in\mathbb{N}$），满足\eqref{align:ch07-7.58}及
\begin{align}
0<U_p'\leqslant U_{p+1}'\leqslant1,\quad \lim_{p\to\infty}U_p'(y)=1\ \text{在 }[0,\infty), \label{align:ch07-7.127}
\end{align}
从而 $\lim_{p\to\infty}U_p(y)=y$（对所有 $y\in[0,\infty)$）。由单调收敛定理、\eqref{align:ch07-7.95}与\eqref{align:ch07-7.109a}得
\begin{align}
\lim_{p\to\infty}U_p^*=\lim_{p\to\infty}\int_{[0,\infty)}U_p(y)\,\mathrm{d}\nu(y)=u_r\quad\text{在 }A\text{ 中几乎处处}. \label{align:ch07-7.128}
\end{align}
我们证明
\begin{align}
\lim_{p\to\infty}[\Lambda_{U_p}^*-\Lambda_{U_p}(U_p^{-1}(U_p^*))]=\phi^*-\phi(u_r)\quad\text{在 }A\text{ 中几乎处处}. \label{align:ch07-7.129}
\end{align}
利用\eqref{align:ch07-7.128}--\eqref{align:ch07-7.129}即可完成证明。事实上，由\eqref{align:ch07-7.128}与控制收敛定理，对几乎所有 $(x,t)\in A$，
\begin{align}
\lim_{p\to\infty}\int_{[0,\infty)}|U_p(y)-U_p^*(x,t)|\,\mathrm{d}\nu_{(x,t)}(y)
=\int_{[0,\infty)}|y-u_r(x,t)|\,\mathrm{d}\nu_{(x,t)}(y). \label{align:ch07-7.130}
\end{align}
在\eqref{align:ch07-7.126}（取 $U=U_p$）中令 $p\to\infty$，由\eqref{align:ch07-7.129}--\eqref{align:ch07-7.130}，对几乎所有 $(x,t)\in A$ 得
\begin{align*}
[\phi^*(x,t)-\phi(u_r)(x,t)]\int_{[0,\infty)}|y-u_r(x,t)|\,\mathrm{d}\nu_{(x,t)}(y)=0.
\end{align*}
由 $A$ 任意，上式蕴含\eqref{align:ch07-7.119}。

余下需证\eqref{align:ch07-7.129}。注意
\begin{align}
\Lambda_{U_p}^*-\Lambda_{U_p}(U_p^{-1}(U_p^*))
=\int_{[0,\infty)}\left(\int_0^y U_p'(z)\phi'(z)\,\mathrm{d}z\right)\mathrm{d}\nu(y)
-\int_0^{U_p^{-1}(U_p^*)}U_p'(z)\phi'(z)\,\mathrm{d}z. \label{align:ch07-7.131}
\end{align}
由\eqref{align:ch07-7.1}与\eqref{align:ch07-7.127}，$U_p'(y)\to1$ 且 $|U_p'(y)\phi'(y)|\leqslant M$（对所有 $y\geqslant0$），由控制收敛定理，
\begin{align}
\lim_{p\to\infty}\Lambda_{U_p}^*(x,t)=\phi^*(x,t)\quad\text{对几乎所有 }(x,t)\in A. \label{align:ch07-7.132}
\end{align}
另一方面，
\begin{align}
\int_0^{U_p^{-1}(U_p^*(x,t))}U_p'(z)\phi'(z)\,\mathrm{d}z-\phi(u_r)(x,t)
=\int_0^{u_r(x,t)}[U_p'(z)-1]\phi'(z)\,\mathrm{d}z
+\int_{u_r(x,t)}^{U_p^{-1}(U_p^*(x,t))}U_p'(z)\phi'(z)\,\mathrm{d}z. \label{align:ch07-7.133}
\end{align}
如\eqref{align:ch07-7.132}那样论证，
\begin{align}
\lim_{p\to\infty}\int_0^{u_r(x,t)}[U_p'(z)-1]\phi'(z)\,\mathrm{d}z=0
\quad\text{对几乎所有 }(x,t)\in A. \label{align:ch07-7.134}
\end{align}
关于\eqref{align:ch07-7.133}右端第二项，对某 $\delta>0$，
\begin{align*}
\left|\int_{u_r(x,t)}^{U_p^{-1}(U_p^*(x,t))}U_p'(z)\phi'(z)\,\mathrm{d}z\right|
\leqslant M|u_r(x,t)-U_p^{-1}(U_p^*(x,t))|
\end{align*}
对充分大的 $p$（见\eqref{align:ch07-7.127}）。由\eqref{align:ch07-7.128}令 $p\to\infty$，得
\begin{align}
\lim_{p\to\infty}\int_{u_r(x,t)}^{U_p^{-1}(U_p^*(x,t))}U_p'(z)\phi'(z)\,\mathrm{d}z=0
\quad\text{对几乎所有 }(x,t)\in A. \label{align:ch07-7.135}
\end{align}
由\eqref{align:ch07-7.133}--\eqref{align:ch07-7.135}得
\begin{align}
\int_0^{U_p^{-1}(U_p^*(x,t))}U_p'(z)\phi'(z)\,\mathrm{d}z\to\phi(u_r)(x,t)
\quad\text{对几乎所有 }(x,t)\in A. \label{align:ch07-7.136}
\end{align}
由\eqref{align:ch07-7.131}、\eqref{align:ch07-7.132}与\eqref{align:ch07-7.136}得\eqref{align:ch07-7.129}。证毕。
\end{proof}

\begin{proof}[\cref{theorem:ch07-7.5.1}的证明]
设 $u\in L^\infty_{w^*}(0,T;\mathscr{R}_f(\Omega))$ 由\eqref{align:ch07-7.108}定义。则由\cref{proposition:ch07-7.5.2}--\cref{proposition:ch07-divcurl}得结论。
\end{proof}

现可证明\cref{theorem:ch07-7.3.1}。

\begin{proof}[\cref{theorem:ch07-7.3.1}的证明]
证明由\eqref{align:ch07-7.108}定义的 $u\in L^\infty_{w^*}(0,T;\mathscr{R}_f(\Omega))$ 满足等式
\begin{align}
\iint_S[u_r\partial_t\zeta+\phi(u_r)\partial_x\zeta]\,\mathrm{d}x\,\mathrm{d}t
+\int_0^T\langle u_s(\cdot,t),\partial_t\zeta(\cdot,t)\rangle\,\mathrm{d}t
=-\langle u_0,\zeta(\cdot,0)\rangle \label{align:ch07-7.137}
\end{align}
对所有 $\zeta\in C^1([0,T];C^1_c(\mathbb{R}))$、$\zeta(\cdot,T)=0$ 在 $\mathbb{R}$ 上成立，从而 $u$ 是 (P) 的解（见\cref{definition:ch07-熵解}）。设 $\{u_{n_k}\}\subseteq\{u_n\}$ 为\cref{theorem:ch07-7.5.1}给出的子列。对 $U(y)=y$ 与 $U(y)=\phi(y)$ 应用\cref{theorem:ch07-7.5.1}(iii)，得
\begin{align}
\lim_{k\to\infty}\iint_S u_{n_k}\partial_t\zeta\,\mathrm{d}x\,\mathrm{d}t
=\iint_S u_r\partial_t\zeta\,\mathrm{d}x\,\mathrm{d}t+\int_0^T\langle u_s(\cdot,t),\partial_t\zeta(\cdot,t)\rangle\,\mathrm{d}t, \label{align:ch07-7.138a}
\end{align}
\begin{align}
\lim_{k\to\infty}\iint_S\phi(u_{n_k})\partial_x\zeta\,\mathrm{d}x\,\mathrm{d}t
=\iint_S\phi(u_r)\partial_x\zeta\,\mathrm{d}x\,\mathrm{d}t \label{align:ch07-7.138b}
\end{align}
（见\eqref{align:ch07-7.89}、\eqref{align:ch07-7.4}与\eqref{align:ch07-7.119}）。由\eqref{align:ch07-7.45}中第一个收敛与\eqref{align:ch07-7.138a}，在\eqref{align:ch07-7.78}（取 $n=n_k$）中令 $k\to\infty$ 得\eqref{align:ch07-7.137}。不等式\eqref{align:ch07-7.5}类似地在\eqref{align:ch07-7.77}（取 $n=n_k$）中令 $k\to\infty$ 证明，对 $U(y)=E_l(y)$ 与 $U(y)=F_l(y)$ 应用\cref{theorem:ch07-7.5.1}(iii)，并如\cref{proposition:ch07-7.5.3}证明中那样论证 $E_l(u_{0n_k})\stackrel{*}{\rightharpoonup}E_l(u_{0r})+u_{0s}$ 在 $\mathscr{R}_f(\mathbb{R})$ 中。结论得证。
\end{proof}

为证\cref{theorem:ch07-7.3.2}，需要下列结果，它刻画 Young 测度 $\nu$ 的离散分解。

\begin{proposition}\label{proposition:ch07-7.5.7}
设 (A0)--(A2) 成立。设 $u\in L^\infty_{w^*}(0,T;\mathscr{R}_f(\mathbb{R}))$ 与 $\nu\in\mathcal{Y}^+(S;\mathbb{R})$ 由\cref{theorem:ch07-7.5.1}给出。则对几乎所有 $(x,t)\in S$，
\begin{align}
\nu_{(x,t)}=\delta_{u_r(x,t)}. \label{align:ch07-7.139}
\end{align}
\end{proposition}

\begin{proof}
固定 $(x,t)\in S$。若 $u_r(x,t)=0$，由\eqref{align:ch07-7.89}立即得\eqref{align:ch07-7.139}。设 $y_1\equiv u_r(x,t)>0$。对所有 $y_2>y_1$ 与充分大的 $p\in\mathbb{N}$，置
\begin{align*}
V_p(y)\triangleq p(y-y_1)\chi_{(y_1,y_1+1/p)}(y)+\chi_{[y_1+1/p,y_2)}(y)+p(y_2+1/p-y)\chi_{[y_2,y_2+1/p)}(y),
\end{align*}
\begin{align*}
\Lambda_{V_p}(y)\triangleq \int_0^y V_p'(z)\phi'(z)\,\mathrm{d}z\quad(y\in[0,\infty)).
\end{align*}
则对任意 $y\in[0,\infty)$，
\begin{align}
\lim_{p\to\infty}V_p(y)=\chi_{(y_1,y_2]}(y), \label{align:ch07-7.140a}
\end{align}
\begin{align}
\lim_{p\to\infty}\Lambda_{V_p}(y)=\phi'(y_1)\chi_{(y_1,y_2]}(y)+[\phi'(y_1)-\phi'(y_2)]\chi_{(y_2,\infty)}(y). \label{align:ch07-7.140b}
\end{align}
设 $\{U_p\}$ 为\cref{proposition:ch07-7.5.6}证明中所用序列。由标准逼近论证，可在等式\eqref{align:ch07-7.125}中取 $U=U_p$、$V=V_p$，得
\begin{align*}
\int_{[0,\infty)}[\Lambda_{U_p}(y)-\Lambda_{U_p}^*(x,t)]V_p(y)\,\mathrm{d}\nu_{(x,t)}(y)
=\int_{[0,\infty)}[U_p(y)-U_p^*(x,t)]\Lambda_{V_p}(y)\,\mathrm{d}\nu_{(x,t)}(y).
\end{align*}
由\eqref{align:ch07-7.128}、\eqref{align:ch07-7.132}与\eqref{align:ch07-7.119}，对所有 $y\geqslant0$ 有
\begin{align*}
\lim_{p\to\infty}U_p(y)-U_p^*(x,t)=y-y_1,
\end{align*}
\begin{align*}
\lim_{p\to\infty}\Lambda_{U_p}^*(x,t)-\Lambda_{U_p}(y)=\phi(y_1)-\phi(y).
\end{align*}
由\eqref{align:ch07-7.140a}--\eqref{align:ch07-7.140b}，在 $S$ 中几乎处处得
\begin{align*}
\lim_{p\to\infty}\int_{[0,\infty)}[\Lambda_{U_p}(y)-\Lambda_{U_p}^*]V_p(y)\,\mathrm{d}\nu(y)
=\int_{(y_1,y_2]}[\phi(y)-\phi(y_1)]\,\mathrm{d}\nu(y),
\end{align*}
\begin{align*}
\begin{aligned}
\lim_{p\to\infty}\int_{[0,\infty)}&[U_p(y)-U_p^*]\Lambda_{V_p}(y)\,\mathrm{d}\nu(y)\\
&=\int_{(y_1,y_2]}\phi'(y_1)(y-y_1)\,\mathrm{d}\nu(y)
+[\phi'(y_1)-\phi'(y_2)]\int_{(y_2,\infty)}(y-y_1)\,\mathrm{d}\nu(y),
\end{aligned}
\end{align*}
从而
\begin{align}
\int_{(y_1,y_2]}[\phi(y)-\phi(y_1)-\phi'(y_1)(y-y_1)]\,\mathrm{d}\nu(y)
=[\phi'(y_1)-\phi'(y_2)]\int_{(y_2,\infty)}(y-y_1)\,\mathrm{d}\nu(y). \label{align:ch07-7.142a}
\end{align}
类似地，对任意 $y_0\in(0,y_1)$，
\begin{align}
\begin{aligned}
\int_{(y_0,y_1)}&[\phi(y)-\phi(y_1)-\phi'(y_1)(y-y_1)]\,\mathrm{d}\nu_{(x,t)}(y)\\
&=[\phi'(y_1)-\phi'(y_0)]\int_{[0,y_0]}(y-y_1)\,\mathrm{d}\nu_{(x,t)}(y).
\end{aligned} \label{align:ch07-7.142b}
\end{align}
由假设 (A2)，$\phi'$ 至少在区间 $[y_0,y_1]$ 与 $[y_1,y_2]$ 之一上严格单调。设在 $[y_1,y_2]$ 上，另一情形类似。则由\eqref{align:ch07-7.142a}，对几乎所有 $(x,t)\in S$，
\begin{align*}
\int_{(y_1,y_2]}|\phi(y)-\phi(y_1)-\phi'(y_1)(y-y_1)|\,\mathrm{d}\nu_{(x,t)}(y)
+|\phi'(y_1)-\phi'(y_2)|\int_{(y_2,\infty)}|y_1-y|\,\mathrm{d}\nu_{(x,t)}(y)=0,
\end{align*}
故 $\operatorname{supp}\nu_{(x,t)}\subseteq[0,y_1]$。因 $\nu_{(x,t)}$ 为概率测度，
\begin{align*}
u_r(x,t)=\int_{[0,u_r(x,t)]}y\,\mathrm{d}\nu_{(x,t)}(y)
=\int_{[0,u_r(x,t)]}[y-u_r(x,t)]\,\mathrm{d}\nu_{(x,t)}(y)+u_r(x,t),
\end{align*}
从而
\begin{align*}
\int_{[0,u_r(x,t)]}|y-u_r(x,t)|\,\mathrm{d}\nu_{(x,t)}(y)=0.
\end{align*}
由上式得 $\operatorname{supp}\nu_{(x,t)}=\{u_r(x,t)\}$，因 $\nu_{(x,t)}$ 为概率测度，\eqref{align:ch07-7.139}得证。
\end{proof}

\begin{remark}\label{remark:ch07-7.5.2}
由\eqref{align:ch07-7.139}与\cref{proposition:ch05-5.4.1}，可能抽取子列（不重新标号），有
\begin{align}
u_{n_k}\to u_r\quad\text{在 }S\text{ 中几乎处处}. \label{align:ch07-7.143}
\end{align}
特别地，若 $\phi$ 有界，由控制收敛定理，对任意 $L>0$，$\phi(u_{n_k})\to\phi(u_r)$ 在 $L^1(R_L)$ 中。
\end{remark}

\begin{remark}\label{remark:ch07-7.5.3}
设 (A0)--(A1) 成立，且对某 $a>0$、$b>0$ 与 $(x,t)\in S$，$\phi'$ 在区间 $I_{a,b}=[u_r(x,t)-a,u_r(x,t)+b]$ 上为常数。此情形可证 $\operatorname{supp}\nu_{(x,t)}\subseteq I_{(x,t)}$，其中 $I_{(x,t)}$ 为 $\phi'=\phi'(u_r)$ 的最大区间（见文献~\cite{ref18} 命题 5.9）。
\end{remark}

\begin{proof}[\cref{theorem:ch07-7.3.2}的证明]
由\cref{theorem:ch07-7.3.1}，$u$ 是 (P) 的 Young 测度熵解，故不等式\eqref{align:ch07-7.5}成立。由\eqref{align:ch07-7.5}与\eqref{align:ch07-7.139}得\eqref{align:ch07-7.8}，结论得证。
\end{proof}

最后证明\cref{lemma:ch07-7.5.4}。

\begin{proof}[\cref{lemma:ch07-7.5.4}的证明]
由标准逼近论证，只需证等式\eqref{align:ch07-7.106a}。

(i) 先在 $U\in C^2([0,\infty))$、$U'\in L^\infty(0,\infty)$ 满足\eqref{align:ch07-7.58}与\eqref{align:ch07-7.91}的假设下证\eqref{align:ch07-7.106a}。设 $\rho\in C^2_c(\mathbb{R})$，$h\in C_c(0,T)$，$\operatorname{supp}\rho\subseteq I_L$（某 $L>0$）。设 $U_{n_k,\rho}$ 由\eqref{align:ch07-7.60}定义。由\eqref{align:ch07-7.104}得
\begin{align}
\lim_{k\to\infty}\int_0^T U_{n_k,\rho}(t)h(t)\,\mathrm{d}t
=\int_0^T U_\rho^*(t)h(t)\,\mathrm{d}t+M_U\int_0^T\langle\sigma(\cdot,t),\rho\rangle h(t)\,\mathrm{d}t. \label{align:ch07-7.144}
\end{align}
另一方面，由\cref{lemma:ch07-7.4.6}，序列 $\{U_{n_k,\rho}\}$ 在 $BV(0,T)$ 中有界，故存在子列（不重新标号）在 $L^1(0,T)$ 中收敛于 $U_\rho^*+M_U\langle\sigma(\cdot,\cdot),\rho\rangle$（见\eqref{align:ch07-7.144}）。因 $\{U(u_{n_k})\}$ 在 $L^\infty(0,T;L^1(I_L))$ 中有界且 $U^*\in L^\infty(0,T;L^1(I_L))$（见\cref{proposition:ch07-7.5.3}与不等式\eqref{align:ch07-7.101}），条件 $\rho\in C^2_c(\mathbb{R})$ 可放宽为 $\rho\in C_c(\mathbb{R})$。故此情形下\eqref{align:ch07-7.106a}成立。

(ii) 现对所有 $U\in C_b([0,\infty))$ 证\eqref{align:ch07-7.106a}（注意此情形 $M_U=0$）。置
\begin{align*}
U_p:[0,\infty)\to\mathbb{R},\quad U_p(y)\triangleq ((U\chi_{[0,p]})*\theta_p)(y)\quad(y\geqslant0),
\end{align*}
其中 $\theta_p(y)\triangleq p\theta(py)$，$\theta\in C^\infty_c(\mathbb{R})$ 为标准磨光核（$p\in\mathbb{N}$）。则 $\{U_p\}\subseteq C^2_c([0,\infty))$，$U_p\to U$ 在 $[0,\infty)$ 的紧子集上一致收敛，且 $\|U_p\|_{L^\infty(\mathbb{R})}\leqslant\|U\|_{L^\infty(\mathbb{R})}$。此外 $U_p'\in L^\infty(0,\infty)$，$U_p$ 满足\eqref{align:ch07-7.58}与\eqref{align:ch07-7.91}（$M_{U_p}=0$），故由 (i) 得
\begin{align}
\lim_{k\to\infty}\int_0^T\left|\int_{\mathbb{R}}U_p(u_{n_k}(x,t))\rho(x)\,\mathrm{d}x-U_{p,\rho}^*(t)\right|\mathrm{d}t=0, \label{align:ch07-7.145}
\end{align}
其中 $U_p^*\triangleq \int_{[0,\infty)}U_p(y)\,\mathrm{d}\nu(y)$，$U_{p,\rho}^*\triangleq \int_{\mathbb{R}}U_p^*(x,t)\rho(x)\,\mathrm{d}x$（对所有 $\rho\in C_c(\mathbb{R})$，$p\in\mathbb{N}$）。现注意
\begin{align}
\begin{aligned}
\int_0^T&\left|\int_{\mathbb{R}}U(u_{n_k}(x,t))\rho(x)\,\mathrm{d}x-U_\rho^*(t)\right|\mathrm{d}t\\
&\leqslant\iint_S|U(u_{n_k}(x,t))-U_p(u_{n_k}(x,t))|\,|\rho(x)|\,\mathrm{d}x\,\mathrm{d}t\\
&\quad+\int_0^T\left|\int_{\mathbb{R}}U_p(u_{n_k}(x,t))\rho(x)\,\mathrm{d}x-U_{p,\rho}^*(t)\right|\mathrm{d}t
+\int_0^T|U_\rho^*(t)-U_{p,\rho}^*(t)|\,\mathrm{d}t.
\end{aligned} \label{align:ch07-7.146}
\end{align}
关于\eqref{align:ch07-7.146}右端第一项，对所有 $p\in\mathbb{N}$ 与 $M>0$，由 Chebyshev 不等式与\eqref{align:ch07-7.71}中第二个不等式，
\begin{align}
\begin{aligned}
\iint_S&|U(u_{n_k})-U_p(u_{n_k})|\,|\rho|\,\mathrm{d}x\,\mathrm{d}t\\
&\leqslant\iint_{\{u_{n_k}\leqslant M\}}|U(u_{n_k})-U_p(u_{n_k})|\,|\rho|\,\mathrm{d}x\,\mathrm{d}t
+\iint_{\{u_{n_k}>M\}}|U(u_{n_k})-U_p(u_{n_k})|\,|\rho|\,\mathrm{d}x\,\mathrm{d}t\\
&\leqslant\|\rho\|_{L^\infty(\mathbb{R})}T\left\{\lambda(\operatorname{supp}\rho)\|U-U_p\|_{L^\infty(0,M)}+\frac{2}{M}\|u_0\|_{\mathscr{R}_f(\mathbb{R})}\|U\|_{L^\infty(\mathbb{R})}\right\}.
\end{aligned} \label{align:ch07-7.147}
\end{align}
关于\eqref{align:ch07-7.146}右端第三项，
\begin{align}
\begin{aligned}
\int_0^T|U_\rho^*(t)-U_{p,\rho}^*(t)|\,\mathrm{d}t
&\leqslant\iint_S|U^*-U_p^*|\,|\rho|\,\mathrm{d}x\,\mathrm{d}t\\
&\leqslant\|\rho\|_{L^\infty(\mathbb{R})}\left\{\lambda(\operatorname{supp}\rho)T\|U-U_p\|_{L^\infty(0,M)}
+2\|U\|_{L^\infty(\mathbb{R})}\iint_{\operatorname{supp}\rho\times(0,T)}\nu_{(x,t)}(\{y>M\})\,\mathrm{d}x\,\mathrm{d}t\right\}.
\end{aligned} \label{align:ch07-7.148}
\end{align}
由\eqref{align:ch07-7.145}--\eqref{align:ch07-7.148}，对所有 $\rho\in C_c(\mathbb{R})$、$p\in\mathbb{N}$ 与 $M>0$，
\begin{align*}
\begin{aligned}
\limsup_{k\to\infty}\int_0^T\mathrm{d}t\,&\left|\int_{\mathbb{R}}U(u_{n_k}(x,t))\rho(x)\,\mathrm{d}x-U_\rho^*(t)\right|\\
&\leqslant2\|\rho\|_{L^\infty(\mathbb{R})}\Big\{\lambda(\operatorname{supp}\rho)T\|U-U_p\|_{L^\infty(0,M)}
+\|U\|_{L^\infty(\mathbb{R})}\left[\frac{T}{M}+\iint_{\operatorname{supp}\rho\times(0,T)}\nu_{(x,t)}(\{y>M\})\,\mathrm{d}x\,\mathrm{d}t\right]\Big\}.
\end{aligned}
\end{align*}
因 $U_p\to U$ 在 $[0,\infty)$ 紧集上一致收敛，令 $p\to\infty$ 得
\begin{align}
\limsup_{k\to\infty}\int_0^T\left|\int_{\mathbb{R}}U(u_{n_k})\rho(x)\,\mathrm{d}x-U_\rho^*(t)\right|\mathrm{d}t
\leqslant2\|\rho\|_{L^\infty(\mathbb{R})}\|U\|_{L^\infty(\mathbb{R})}\left\{\frac{T}{M}+\iint_{\operatorname{supp}\rho\times(0,T)}\nu_{(x,t)}(\{y>M\})\,\mathrm{d}x\,\mathrm{d}t\right\}. \label{align:ch07-7.149}
\end{align}
因 $\nu_{(x,t)}$ 为概率测度，由控制收敛定理显然
\begin{align*}
\lim_{M\to\infty}\iint_{\operatorname{supp}\rho\times(0,T)}\nu_{(x,t)}(\{y>M\})\,\mathrm{d}x\,\mathrm{d}t=0.
\end{align*}
在\eqref{align:ch07-7.149}中令 $M\to\infty$，即得\eqref{align:ch07-7.106a}。

(iii) 现设 $U\in C([0,\infty))$ 为任意满足\eqref{align:ch07-7.91}的函数。如\cref{proposition:ch07-7.5.3}证明中那样，设 $f_{j1},f_{j2}\in C^2([0,\infty))$ 满足\eqref{align:ch07-7.100}（$j\in\mathbb{N}$，$j\geqslant[y_\epsilon]+1$）。则
\begin{align*}
U=U_{j1}+U_{j2},\quad U^*=U_{j1}^*+U_{j2}^*,\quad U_\rho^*=U_{j1,\rho}^*+U_{j2,\rho}^*,
\end{align*}
其中 $U_{jl}\triangleq f_{jl}U$，$U_{jl}^*\triangleq \int_{[0,\infty)}U_{jl}(y)\,\mathrm{d}\nu(y)$，$U_{jl,\rho}^*\triangleq \int_{\mathbb{R}}U_{jl}^*(x,t)\rho(x)\,\mathrm{d}x$（$\rho\in C_c(\mathbb{R})$，$l=1,2$）。因此
\begin{align}
\begin{aligned}
\int_0^T&\left|\int_{\mathbb{R}}U(u_{n_k}(x,t))\rho(x)\,\mathrm{d}x-U_\rho^*(t)-M_U\langle\sigma(\cdot,t),\rho\rangle\right|\mathrm{d}t\\
&\leqslant\int_0^T\left|\int_{\mathbb{R}}U_{j1}(u_{n_k}(x,t))\rho(x)\,\mathrm{d}x-U_{j1,\rho}^*(t)\right|\mathrm{d}t\\
&\quad+\int_0^T\left|\int_{\mathbb{R}}U_{j2}(u_{n_k}(x,t))\rho(x)\,\mathrm{d}x-U_{j2,\rho}^*(t)-M_U\langle\sigma(\cdot,t),\rho\rangle\right|\mathrm{d}t=:I_{jk1}+I_{jk2}.
\end{aligned} \label{align:ch07-7.151}
\end{align}
因 $U_{j1}\in C_b([0,\infty))$ 从而 $M_{U_{j1}}=0$，由 (ii) 与\eqref{align:ch07-7.106a}得
\begin{align}
\lim_{k\to\infty}I_{jk1}=0\quad(j\in\mathbb{N},\ j\geqslant[y_\epsilon]+1). \label{align:ch07-7.152}
\end{align}
另一方面，
\begin{align}
\begin{aligned}
I_{jk2}\leqslant{}&\int_0^T\left|\int_{\mathbb{R}}[U_{j2}(u_{n_k}(x,t))-M_U f_{j2}(u_{n_k}(x,t))u_{n_k}(x,t)]\rho(x)\,\mathrm{d}x\right|\mathrm{d}t\\
&+M_U\int_0^T\mathrm{d}t\left|\int_{\mathbb{R}}f_{j2}(u_{n_k}(x,t))u_{n_k}(x,t)\rho(x)\,\mathrm{d}x-\int_{\mathbb{R}}F_{j2}^*(x,t)\rho(x)\,\mathrm{d}x-\langle\sigma(\cdot,t),\rho\rangle\right|\\
&+\iint_S|U_{j2}^*-M_U F_{j2}^*|(x,t)\,|\rho(x)|\,\mathrm{d}x\,\mathrm{d}t,
\end{aligned} \label{align:ch07-7.153}
\end{align}
其中 $F_{j2}^*\triangleq \int_{[0,\infty)}f_{j2}(y)y\,\mathrm{d}\nu(y)$。

分别处理\eqref{align:ch07-7.153}右端三项。由\eqref{align:ch07-7.99}与\eqref{align:ch07-7.71}中第二个不等式，
\begin{align}
\int_0^T\left|\int_{\mathbb{R}}[U_{j2}(u_{n_k}(x,t))-M_U f_{j2}(u_{n_k}(x,t))u_{n_k}(x,t)]\rho(x)\,\mathrm{d}x\right|\mathrm{d}t
\leqslant\epsilon\|\rho\|_{L^\infty(\mathbb{R})}T\|u_0\|_{\mathscr{R}_f(\mathbb{R})}. \label{align:ch07-7.154}
\end{align}
此外，映射 $y\mapsto f_{j2}(y)y$ 属于 $C^2([0,\infty))$，导数有界，满足\eqref{align:ch07-7.58}与\eqref{align:ch07-7.91}（$M_U=1$）。故由 (i)，
\begin{align}
\lim_{k\to\infty}\int_0^T\mathrm{d}t\left|\int_{\mathbb{R}}f_{j2}(u_{n_k}(x,t))u_{n_k}(x,t)\rho(x)\,\mathrm{d}x-\int_{\mathbb{R}}F_{j2}^*(x,t)\rho(x)\,\mathrm{d}x-\langle\sigma(\cdot,t),\rho\rangle\right|=0. \label{align:ch07-7.155}
\end{align}
现证
\begin{align}
\lim_{j\to\infty}\iint_S|U_{j2}^*-M_U F_{j2}^*|(x,t)\,|\rho(x)|\,\mathrm{d}x\,\mathrm{d}t=0. \label{align:ch07-7.156}
\end{align}
对几乎所有 $(x,t)\in S$，
\begin{align*}
|U_{j2}^*(x,t)-M_U F_{j2}^*(x,t)|
\leqslant\int_{[0,\infty)}|U_{j2}(y)-M_U f_{j2}(y)y|\,\mathrm{d}\nu_{(x,t)}(y)
=\int_{[0,\infty)}f_{j2}(y)|U(y)-M_U y|\,\mathrm{d}\nu_{(x,t)}(y)
\leqslant\epsilon\int_{[j,\infty)}y\,\mathrm{d}\nu_{(x,t)}(y)\leqslant\epsilon u_b(x,t)
\end{align*}
（见\eqref{align:ch07-7.95}）。因 $u_b\in L^\infty(0,T;L^1(\mathbb{R}))$ 且对几乎所有 $(x,t)\in S$ 有 $\lim_{j\to\infty}\int_{[j,\infty)}y\,\mathrm{d}\nu_{(x,t)}(y)=0$，由控制收敛定理得等式\eqref{align:ch07-7.156}。

由\eqref{align:ch07-7.153}--\eqref{align:ch07-7.156}，
\begin{align*}
\limsup_{j\to\infty}\left(\limsup_{k\to\infty}I_{jk2}\right)\leqslant\epsilon\|\rho\|_{L^\infty(\mathbb{R})}T\|u_0\|_{\mathscr{R}_f(\mathbb{R})},
\end{align*}
故由\eqref{align:ch07-7.151}--\eqref{align:ch07-7.152}，
\begin{align*}
\limsup_{k\to\infty}\int_0^T\left|\int_{\mathbb{R}}[U(u_{n_k})-U^*](x,t)\rho(x)\,\mathrm{d}x-M_U\langle\sigma(\cdot,t),\rho\rangle\right|\mathrm{d}t
\leqslant\epsilon T\|u_0\|_{\mathscr{R}_f(\mathbb{R})}\|\rho\|_{L^\infty(\mathbb{R})}.
\end{align*}
由 $\epsilon$ 任意，上式给出\eqref{align:ch07-7.106a}。证毕。
\end{proof}




\end{document}